\documentclass[lang=en,11pt]{elegantbook}

% ── Packages ──
\usepackage{booktabs}
\usepackage{longtable}
\usepackage{tabularx}
\usepackage{listings}
\usepackage{xcolor}
\usepackage{caption}
\usepackage{enumitem}

% ── Colors ──
\definecolor{slipkw}{RGB}{0,0,180}
\definecolor{sliptick}{RGB}{0,128,128}
\definecolor{svkw}{RGB}{0,0,180}
\definecolor{codegreen}{RGB}{80,120,80}
\definecolor{codegray}{RGB}{120,120,120}
\definecolor{codebg}{RGB}{248,248,248}

% ── Slip language definition ──
\lstdefinelanguage{Slip}{
  keywords={module,param,localparam,logic,signed,assign,seq,comb,initial,
            pos,neg,if,else,for,case,casez,casex,default,inside,begin,end,include},
  sensitive=true,
  morecomment=[l]{//},
  morecomment=[s]{/*}{*/},
  morestring=[b]",
  literate={`for}{{\textcolor{sliptick}{\textbf{`for}}}}{4}
           {`if}{{\textcolor{sliptick}{\textbf{`if}}}}{3}
           {`else}{{\textcolor{sliptick}{\textbf{`else}}}}{5}
           {`i}{{\textcolor{sliptick}{\textbf{`i}}}}{2}
           {`j}{{\textcolor{sliptick}{\textbf{`j}}}}{2},
}

% ── SystemVerilog language definition ──
\lstdefinelanguage{SystemVerilog}{
  keywords={module,endmodule,input,output,inout,logic,wire,reg,
            always_ff,always_comb,always_latch,begin,end,
            if,else,for,parameter,localparam,assign,
            posedge,negedge,generate,endgenerate,genvar,
            case,endcase,default,initial,integer},
  sensitive=true,
  morecomment=[l]{//},
  morecomment=[s]{/*}{*/},
  morestring=[b]",
}

% ── Listings global style ──
\lstset{
  basicstyle=\ttfamily\footnotesize,
  keywordstyle=\color{slipkw}\bfseries,
  keywordstyle=[2]\color{sliptick}\bfseries,
  commentstyle=\color{codegray}\itshape,
  stringstyle=\color{codegreen},
  backgroundcolor=\color{codebg},
  numbers=none,
  frame=single,
  rulecolor=\color{gray!40},
  framerule=0.4pt,
  breaklines=true,
  showstringspaces=false,
  tabsize=4,
  xleftmargin=2pt,
  xrightmargin=2pt,
  aboveskip=8pt,
  belowskip=4pt,
}

% ── Inline code shorthand ──
\newcommand{\code}[1]{\texttt{#1}}
\newcommand{\slipinline}[1]{\lstinline[language=Slip]|#1|}
\newcommand{\svinline}[1]{\lstinline[language=SystemVerilog]|#1|}

% ── Disable first-line indent (English typesetting) ──
\setlength{\parindent}{0pt}
\setlength{\parskip}{6pt}

% ── Title metadata ──
\title{Slip User Manual}
\subtitle{A Minimal HDL DSL Compiler}
\author{Slip Project}
\date{May 2026}
\version{1.0}

\begin{document}

\maketitle
\frontmatter
\tableofcontents
\mainmatter

% ════════════════════════════════════════════════════════════════════
% Part I: Getting Started
% ════════════════════════════════════════════════════════════════════

\chapter{Introduction}

Slip is a minimal hardware description language (DSL) that compiles \code{.slip} source files into synthesizable SystemVerilog. It is designed as syntactic sugar over SystemVerilog, providing a streamlined syntax for common hardware design patterns while generating clean, human-readable, IEEE 1800-2017 compliant output.

\section{Design Philosophy}

\begin{itemize}
  \item \textbf{Minimal syntax} --- port directions, signal types, and instance wiring are inferred from context.
  \item \textbf{Direct mapping} --- every Slip construct maps 1:1 to a SystemVerilog equivalent.
  \item \textbf{Compile-time metaprogramming} --- \slipinline{`for} loops and \slipinline{`if} conditionals are fully expanded before code generation, with automatic constant folding.
  \item \textbf{Validated output} --- generated SystemVerilog is checked against the IEEE 1800-2017 standard via \href{https://github.com/MikePopoloski/slang}{pyslang}.
\end{itemize}

% ──────────────────────────────────────────────────────────────────
\chapter{Installation}
% ──────────────────────────────────────────────────────────────────

\section{Prerequisites}

\begin{itemize}
  \item Python $\geq$ 3.10
  \item \href{https://github.com/astral-sh/uv}{uv} (recommended) or pip
\end{itemize}

\section{Install with uv}

\begin{lstlisting}[language=bash]
git clone https://github.com/user/slip.git
cd slip
uv sync
\end{lstlisting}

\section{Install with pip}

\begin{lstlisting}[language=bash]
git clone https://github.com/user/slip.git
cd slip
pip install -e .
\end{lstlisting}

% ──────────────────────────────────────────────────────────────────
\chapter{Command Line Interface}
% ──────────────────────────────────────────────────────────────────

\section{\code{slip build} --- Compile to SystemVerilog}

\begin{lstlisting}[language=bash]
slip build <file.slip>                        # output to ./build/
slip build -o <out_dir> <file.slip>           # custom output directory
slip build -ip <ip_dir> <file.slip>           # include external IP directory (recursive)
slip build -ip <dir1> -ip <dir2> <file.slip>  # multiple IP directories
slip build -f <file.f> <file.slip>            # use VCS-format filelist
slip build -f <file.f> -ip <dir> <file.slip>  # filelist + directory (combined)
\end{lstlisting}

Each module in the \code{.slip} file generates a separate \code{.sv} file in the output directory.

\section{\code{slip check} --- Validate without generating output}

\begin{lstlisting}[language=bash]
slip check <file.slip>                        # validate only
slip check -ip <ip_dir> <file.slip>           # with IP directory
slip check -f <file.f> <file.slip>            # with filelist
\end{lstlisting}

Runs the full compilation pipeline (parsing, semantic analysis, code generation) but discards the output. Useful for CI validation.

% ════════════════════════════════════════════════════════════════════
% Part II: Language Reference
% ════════════════════════════════════════════════════════════════════

\chapter{Language Overview}

A Slip source file contains one or more module declarations. Comments use \code{//} (line) or \code{/* */} (block) syntax.

\begin{lstlisting}[language=Slip]
// This is a comment
module AND_gate (a, b, y) {
    logic a;
    logic b;
    logic y;
    assign y = a & b;
}
\end{lstlisting}

Generates:

\begin{lstlisting}[language=SystemVerilog]
module AND_gate (
    input logic a,
    input logic b,
    output logic y
);

    assign y = a & b;

endmodule
\end{lstlisting}

\section{Key Differences from SystemVerilog}

\begin{itemize}
  \item No \svinline{input}/\svinline{output} keywords --- port directions are inferred.
  \item No \svinline{wire}/\svinline{reg} --- everything is \svinline{logic}.
  \item Blocking assignment in \code{seq} blocks is automatically converted to non-blocking.
  \item Compile-time metaprogramming via \slipinline{`for} and \slipinline{`if}.
  \item Template identifiers embed loop variables in signal names (e.g., \code{data\_`i}).
  \item \code{include "file.slip"} directive for reusing \code{defun} functions and modules across files.
  \item Generated SV annotated with \code{// implicit}, \code{// no width}, \code{// width from X.Y} for easy review.
\end{itemize}

% ──────────────────────────────────────────────────────────────────
\chapter{Include Directive}
% ──────────────────────────────────────────────────────────────────

The \code{include} directive imports \code{defun} function definitions and module declarations from another \code{.slip} file. It must appear at the top level, outside module bodies.

\section{Basic Usage}

\begin{lstlisting}[language=Slip]
include "helpers.slip";
include "utils/common.slip";
\end{lstlisting}

Included files are lexed, parsed, and their contents merged into the compilation unit. Relative paths are resolved from the directory of the including file.

\section{Example: Sharing defun Libraries}

File \code{helpers.slip}:

\begin{lstlisting}[language=Slip]
defun prefix(s):
    return "bus_" + s

defun inc(n):
    return n + 1
\end{lstlisting}

File \code{design.slip}:

\begin{lstlisting}[language=Slip]
include "helpers.slip";

module top (clk) {
    logic clk;
    inner u1 {
        .clk,
        "data_(.*)" => "$prefix(\\1)",
        "ch_(.*)" => "$inc(\\1)"
    };
}
\end{lstlisting}

\section{Recursive Includes}

Include directives in included files are also resolved. The same file is never included twice (deduplicated by absolute path).

\section{Restrictions}

\begin{itemize}
  \item \code{include} must appear at the top level --- not inside module bodies.
  \item The included path must be a string literal. Variable-based paths are not supported.
  \item Circular includes are detected and silently skipped.
\end{itemize}

% ──────────────────────────────────────────────────────────────────
\chapter{Module Declaration}
% ──────────────────────────────────────────────────────────────────

\section{Basic Syntax}

\begin{lstlisting}[language=Slip]
module <name> (port1, port2, ...) {
    // body
}
\end{lstlisting}

\section{With Parameters}

\begin{lstlisting}[language=Slip]
module <name> #(param W = 8, param DEPTH = 16) (port1, port2, ...) {
    // body
}
\end{lstlisting}

\section{Implicit Ports}

If the port list is omitted entirely, all identifiers used in the module body that are not declared internally become ports:

\begin{lstlisting}[language=Slip]
module passthrough {
    assign out = in;
}
\end{lstlisting}

Generates:

\begin{lstlisting}[language=SystemVerilog]
module passthrough (
    input logic in,
    output logic out
);

    assign out = in;

endmodule
\end{lstlisting}

Port directions are inferred from usage (see Chapter~\ref{ch:port-inference}).

% ──────────────────────────────────────────────────────────────────
\chapter{Parameters and Localparams}
% ──────────────────────────────────────────────────────────────────

\section{Parameters (\code{param})}

Parameters are declared in the module header and can be overridden during instantiation:

\begin{lstlisting}[language=Slip]
module counter #(param WIDTH = 8, param MAX = 255) (clk, rst_n, count) {
    logic clk;
    logic rst_n;
    logic [WIDTH-1:0] count;

    seq (clk, neg: rst_n) {
        if (!rst_n) {
            count <= 0;
        } else if (count == MAX) {
            count <= 0;
        } else {
            count <= count + 1;
        }
    }
}
\end{lstlisting}

Generates:

\begin{lstlisting}[language=SystemVerilog]
module counter #(
    parameter WIDTH = 8,
    parameter MAX = 255
) (
    input logic clk,
    input logic rst_n,
    output logic [WIDTH - 1:0] count
);

    always_ff @(posedge clk or negedge rst_n) begin
        if (!rst_n) begin
            count <= 0;
        end else begin
            if (count == MAX) begin
                count <= 0;
            end else begin
                count <= count + 1;
            end
        end
    end

endmodule
\end{lstlisting}

Override at instantiation:

\begin{lstlisting}[language=Slip]
module top (clk, rst_n, count) {
    logic clk;
    logic rst_n;
    logic [15:0] count;
    counter #(.WIDTH(16), .MAX(65535)) u_cnt { .clk, .rst_n, .count };
}
\end{lstlisting}

Generates:

\begin{lstlisting}[language=SystemVerilog]
module top (
    input logic clk,
    input logic rst_n,
    input logic [15:0] count
);

    counter #(
        .WIDTH(16),
        .MAX(65535)
    ) u_cnt (
        .clk(clk),
        .rst_n(rst_n),
        .count(count)
    );

endmodule
\end{lstlisting}

\section{Localparams (\code{localparam})}

Localparams are declared in the module body and \textbf{cannot} be overridden during instantiation:

\begin{lstlisting}[language=Slip]
module fsm (clk, state) {
    logic clk;
    logic [1:0] state;
    localparam IDLE = 0;
    localparam RUN  = 1;
    localparam DONE = 2;
}
\end{lstlisting}

Generates:

\begin{lstlisting}[language=SystemVerilog]
module fsm (
    input logic clk,
    input logic [1:0] state
);

    localparam IDLE = 0;
    localparam RUN = 1;
    localparam DONE = 2;

endmodule
\end{lstlisting}

Attempting to override a localparam in an instance raises a semantic error:

\begin{lstlisting}[language=Slip]
fsm #(.IDLE(3)) u1 { .clk, .state };  // ERROR
\end{lstlisting}

Localparams can be used in signal widths, \slipinline{`for} bounds, and \slipinline{`if} conditions.

% ──────────────────────────────────────────────────────────────────
\chapter{Signals}
% ──────────────────────────────────────────────────────────────────

\section{Explicit Declaration}

\begin{lstlisting}[language=Slip]
logic [7:0] data;          // 8-bit signal
logic flag;                // 1-bit signal
signed [15:0] val;         // signed 16-bit
logic [7:0] mem [0:15];    // array of 16 x 8-bit
logic [WIDTH-1:0] bus;     // parameterized width
\end{lstlisting}

The \code{logic} keyword is optional for bare declarations:

\begin{lstlisting}[language=Slip]
data;                      // equivalent to: logic data;
\end{lstlisting}

\section{Implicit Declaration}

Signals that are referenced but not declared are automatically created as \code{logic} signals:

\begin{lstlisting}[language=Slip]
module example (a, b) {
    logic a;
    logic b;
    assign y = a & b;  // 'y' is implicitly declared as logic
}
\end{lstlisting}

Generates:

\begin{lstlisting}[language=SystemVerilog]
module example (
    input logic a,
    input logic b
);

    logic y;

    assign y = a & b;

endmodule
\end{lstlisting}

Implicit signals used in instance port connections automatically inherit the target port's width, annotated with \code{// width from X.Y} in the generated SV. Other implicit signals default to 1-bit and are annotated with \code{// implicit, no width} for easy review.

% ──────────────────────────────────────────────────────────────────
\chapter{Port Inference}
\label{ch:port-inference}
% ──────────────────────────────────────────────────────────────────

Slip infers port directions automatically from signal usage within the module body:

\begin{table}[htbp]
\centering
\begin{tabular}{ll}
\toprule
\textbf{Usage Pattern} & \textbf{Inferred Direction} \\
\midrule
Only appears in expressions (right-hand side) & \code{input} \\
Only appears as assignment target (left-hand side) & \code{output} \\
Both read and driven & \code{inout} \\
\bottomrule
\end{tabular}
\end{table}

\section{Example}

\begin{lstlisting}[language=Slip]
module adder (a, b, sum) {
    logic a;
    logic b;
    logic sum;
    assign sum = a + b;
}
\end{lstlisting}

Generates:

\begin{lstlisting}[language=SystemVerilog]
module adder (
    input logic a,
    input logic b,
    output logic sum
);

    assign sum = a + b;

endmodule
\end{lstlisting}

% ──────────────────────────────────────────────────────────────────
\chapter{Assignments}
% ──────────────────────────────────────────────────────────────────

\section{Continuous Assignment}

At module level, both forms generate \svinline{assign} statements in SystemVerilog:

\begin{lstlisting}[language=Slip]
assign y = a & b;          // explicit 'assign' keyword
y = a + b;                 // bare assignment (also generates 'assign')
\end{lstlisting}

Both generate continuous assignments:

\begin{lstlisting}[language=SystemVerilog]
assign y = a & b;
assign z = a + b;
\end{lstlisting}

\section{Procedural Assignment}

Inside \code{seq} and \code{comb} blocks, assignments become procedural:

\begin{lstlisting}[language=Slip]
seq (clk) {
    q <= d;                // non-blocking (sequential)
}

comb {
    y = sel ? a : b;       // blocking (combinational)
}
\end{lstlisting}

Generates:

\begin{lstlisting}[language=SystemVerilog]
    always_ff @(posedge clk) begin
        q <= d;
    end

    always_comb begin
        y = sel ? a : b;
    end
\end{lstlisting}

\begin{remark}
The \code{assign} keyword must \textbf{not} be used inside \code{seq} or \code{comb} blocks. It is reserved for module-level continuous assignments only.
\end{remark}

% ──────────────────────────────────────────────────────────────────
\chapter{Sequential Logic (\code{seq})}
% ──────────────────────────────────────────────────────────────────

The \code{seq} block generates an \svinline{always\_ff} block with the specified clock and optional asynchronous reset.

\section{Basic Clock}

\begin{lstlisting}[language=Slip]
seq (clk) {
    q <= d;
}
\end{lstlisting}

Generates:

\begin{lstlisting}[language=SystemVerilog]
    always_ff @(posedge clk) begin
        q <= d;
    end
\end{lstlisting}

\section{Clock with Asynchronous Reset}

\begin{lstlisting}[language=Slip]
seq (clk, neg: rst_n) {
    if (!rst_n) {
        q <= 0;
    } else {
        q <= d;
    }
}
\end{lstlisting}

Generates:

\begin{lstlisting}[language=SystemVerilog]
    always_ff @(posedge clk or negedge rst_n) begin
        if (!rst_n) begin
            q <= 0;
        end else begin
            q <= d;
        end
    end
\end{lstlisting}

Reset polarity options:
\begin{itemize}
  \item \code{neg: rst\_n} --- active-low reset (\svinline{negedge rst\_n})
  \item \code{pos: rst} --- active-high reset (\svinline{posedge rst})
\end{itemize}

\section{Automatic Blocking-to-Nonblocking Conversion}

All blocking assignments (\code{=}) inside \code{seq} blocks are automatically converted to non-blocking (\code{<=}). This prevents common sequential logic mistakes:

\begin{lstlisting}[language=Slip]
seq (clk) {
    q = d;   // automatically becomes q <= d;
}
\end{lstlisting}

Generates (blocking \code{=} converted to non-blocking \code{<=}):

\begin{lstlisting}[language=SystemVerilog]
    always_ff @(posedge clk) begin
        q <= d;
    end
\end{lstlisting}

A warning is emitted when this conversion occurs.

% ──────────────────────────────────────────────────────────────────
\chapter{Combinational Logic (\code{comb})}
% ──────────────────────────────────────────────────────────────────

The \code{comb} block generates an \svinline{always\_comb} block. Assignments inside remain blocking:

\begin{lstlisting}[language=Slip]
comb {
    if (sel) {
        y = a;
    } else {
        y = b;
    }
}
\end{lstlisting}

Generates:

\begin{lstlisting}[language=SystemVerilog]
    always_comb begin
        if (sel) begin
            y = a;
        end else begin
            y = b;
        end
    end
\end{lstlisting}

If the combinational block infers a latch (e.g., incomplete assignment in an \code{if} without \code{else}), Slip emits \svinline{always\_latch} without a sensitivity list:

\begin{lstlisting}[language=SystemVerilog]
always_latch begin
    if (en) begin
        q <= d;
    end
end
\end{lstlisting}

% ──────────────────────────────────────────────────────────────────
\chapter{Initial Blocks}
% ──────────────────────────────────────────────────────────────────

The \code{initial} block generates an \svinline{initial} block. It is primarily used in testbenches for initialization sequences. Assignments inside remain blocking:

\begin{lstlisting}[language=Slip]
initial {
    clk = 0;
    rst_n = 0;
}
\end{lstlisting}

Generates:

\begin{lstlisting}[language=SystemVerilog]
    initial begin
        clk = 0;
        rst_n = 0;
    end
\end{lstlisting}

% ──────────────────────────────────────────────────────────────────
\chapter{Case Statements}
% ──────────────────────────────────────────────────────────────────

Slip supports \code{case}, \code{casez}, and \code{casex} statements with a brace-based syntax.

\section{Basic Case}

\begin{lstlisting}[language=Slip]
case (sel) {
    2'b00: y = a;
    2'b01: y = b;
    2'b10, 2'b11: y = c;
    default: y = 0;
}
\end{lstlisting}

Generates:

\begin{lstlisting}[language=SystemVerilog]
    always_comb begin
        case (sel)
            2'b00: begin
                y = a;
            end
            2'b01: begin
                y = b;
            end
            2'b10, 2'b11: begin
                y = c;
            end
            default: begin
                y = 0;
            end
        endcase
    end
\end{lstlisting}

\section{Casez with Wildcards}

\code{casez} treats \code{?} as don't-care bits:

\begin{lstlisting}[language=Slip]
casez (opcode) {
    4'b1???: y = a;
    4'b01??: y = b;
    4'b001?: y = c;
    4'b0001: y = d;
    default: y = 0;
}
\end{lstlisting}

Generates:

\begin{lstlisting}[language=SystemVerilog]
    always_comb begin
        casez (opcode)
            4'b1???: begin
                y = a;
            end
            4'b01??: begin
                y = b;
            end
            4'b001?: begin
                y = c;
            end
            4'b0001: begin
                y = d;
            end
            default: begin
                y = 0;
            end
        endcase
    end
\end{lstlisting}

\section{Multi-Pattern Items}

A single case item can match multiple patterns, separated by commas:

\begin{lstlisting}[language=Slip]
localparam IDLE = 0;
localparam WAIT = 1;
localparam RUN  = 2;

comb {
    case (state) {
        IDLE, WAIT: y = 0;
        RUN: y = 1;
        default: y = 0;
    }
}
\end{lstlisting}

Generates (assuming \code{IDLE=0}, \code{WAIT=1}, \code{RUN=2}):

\begin{lstlisting}[language=SystemVerilog]
localparam IDLE = 0;
localparam WAIT = 1;
localparam RUN = 2;

    always_comb begin
        case (state)
            IDLE, WAIT: begin
                y = 0;
            end
            RUN: begin
                y = 1;
            end
            default: begin
                y = 0;
            end
        endcase
    end
\end{lstlisting}

% ──────────────────────────────────────────────────────────────────
\chapter{Inside Operator}
% ──────────────────────────────────────────────────────────────────

The \code{inside} operator tests set membership, following SystemVerilog 2009:

\begin{lstlisting}[language=Slip]
if (state inside {IDLE, RUN, DONE}) begin
    // state is one of IDLE, RUN, DONE
end
\end{lstlisting}

The right-hand side is a set literal \code{\{v1, v2, ...\}}. At compile time, if all values are constants, the expression evaluates to 0 or 1.

\begin{lstlisting}[language=Slip]
localparam IDLE = 0;
localparam RUN  = 1;
localparam DONE = 2;

comb {
    if (state inside {IDLE, RUN, DONE}) {
        active = 1;
    } else {
        active = 0;
    }
}
\end{lstlisting}

Generates:

\begin{lstlisting}[language=SystemVerilog]
    localparam IDLE = 0;
    localparam RUN = 1;
    localparam DONE = 2;

    logic active;

    always_comb begin
        if (state inside {IDLE, RUN, DONE}) begin
            active = 1;
        end else begin
            active = 0;
        end
    end
\end{lstlisting}

% ──────────────────────────────────────────────────────────────────
\chapter{Compound Assignments}
% ──────────────────────────────────────────────────────────────────

Slip supports all SystemVerilog compound assignment operators. They are desugared at parse time:

\begin{table}[htbp]
\centering
\begin{tabular}{lll}
\toprule
\textbf{Operator} & \textbf{Desugars To} & \textbf{Example} \\
\midrule
\code{+=}  & \code{a = a + b}   & \code{i += 1} \\
\code{-=}  & \code{a = a - b}   & \code{count -= 1} \\
\code{*=}  & \code{a = a * b}   & \code{val *= 2} \\
\code{/=}  & \code{a = a / b}   & \code{val /= 2} \\
\code{\%=} & \code{a = a \% b}  & \code{val \%= N} \\
\code{\&=} & \code{a = a \& b}  & \code{flags \&= MASK} \\
\code{|=}  & \code{a = a | b}   & \code{flags |= BIT} \\
\code{\^=} & \code{a = a \^ b}  & \code{flags \^= BIT} \\
\code{<<=} & \code{a = a << b}  & \code{val <<= 2} \\
\code{>>=} & \code{a = a >> b}  & \code{val >>= 2} \\
\code{<<<=} & \code{a = a <<< b} & \code{val <<<= 1} \\
\code{>>>=} & \code{a = a >>> b} & \code{val >>>= 1} \\
\bottomrule
\end{tabular}
\end{table}

Compound assignments are valid in all contexts where regular assignments are used (module-level, \code{comb}, \code{seq}, \code{initial}).

\section{Example}

\begin{lstlisting}[language=Slip]
module m (clk, val, delta) {
    logic clk;
    logic [7:0] val;
    logic [7:0] delta;

    seq (clk) {
        val += delta;
        val -= 1;
    }
}
\end{lstlisting}

Generates (compound assignments desugared, blocking converted to non-blocking):

\begin{lstlisting}[language=SystemVerilog]
module m (
    input logic clk,
    output logic [7:0] val,
    input logic [7:0] delta
);

    always_ff @(posedge clk) begin
        val <= val + delta;
        val <= val - 1;
    end

endmodule
\end{lstlisting}

% ──────────────────────────────────────────────────────────────────
\chapter{Control Flow}
% ──────────────────────────────────────────────────────────────────

\section{If / Else}

\begin{lstlisting}[language=Slip]
if (condition) begin
    // ...
end else begin
    // ...
end
\end{lstlisting}

\section{For Loop}

Procedural \code{for} loops generate SystemVerilog \svinline{for} statements. They are valid inside \code{seq} and \code{comb} blocks:

\begin{lstlisting}[language=Slip]
seq (clk, neg: rst_n) {
    if (!rst_n) {
        for (i = 0; i < 4; i = i + 1) {
            reg_file[i] <= 0;
        }
    }
}
\end{lstlisting}

Generates:

\begin{lstlisting}[language=SystemVerilog]
    logic [7:0] reg_file [0:3];

    always_ff @(posedge clk or negedge rst_n) begin
        if (!rst_n) begin
            for (int i = 0; i < 4; i = i + 1) begin
                reg_file[i] <= 0;
            end
        end
    end
\end{lstlisting}

Loop variables (e.g., \code{i}) are excluded from implicit signal creation --- they are treated as loop-local integers.

\section{For-Loop Step Variants}

The step clause supports three forms:

\begin{lstlisting}[language=Slip]
for (i = 0; i < N; i = i + 1) { ... }   // standard
for (i = 0; i < N; i += 1)    { ... }   // compound assignment
for (i = 0; i < N; i++)       { ... }   // increment
\end{lstlisting}

Compound assignments (\code{+=}, \code{-=}, \code{*=}, etc.) are desugared at parse time: \code{i += 1} becomes \code{i = i + 1}. All three forms generate the same SystemVerilog output:

\begin{lstlisting}[language=SystemVerilog]
for (int i = 0; i < N; i = i + 1) begin
    // ...
end
\end{lstlisting}

\section{Restrictions}

\begin{itemize}
  \item \code{while} loops are not supported (combinational while loops are unsafe in synthesis).
\end{itemize}

% ──────────────────────────────────────────────────────────────────
\chapter{Module Instantiation}
% ──────────────────────────────────────────────────────────────────

\section{Basic Instantiation}

\begin{lstlisting}[language=Slip]
module_name instance_name {
    .port1(signal1),
    .port2(signal2)
};
\end{lstlisting}

\section{Same-Name Shorthand}

When the port name matches the signal name, the signal expression can be omitted:

\begin{lstlisting}[language=Slip]
inner u1 { .clk, .rst_n, .data_in };
// Equivalent to: .clk(clk), .rst_n(rst_n), .data_in(data_in)
\end{lstlisting}

\section{Parameter Override}

\begin{lstlisting}[language=Slip]
inner #(.WIDTH(16), .DEPTH(32)) u1 { .clk, .data_in, .data_out };
\end{lstlisting}

\section{Dangling Ports (\code{\_})}

Use \code{\_} to leave a port intentionally unconnected:

\begin{lstlisting}[language=Slip]
inner u1 { .clk, .rst_n(_), .data_in };
\end{lstlisting}

In the generated SystemVerilog, the \code{.rst\_n()} connection is omitted entirely.

\section{Regex Port Mapping}

Use regex patterns to map groups of ports to signals:

\begin{lstlisting}[language=Slip]
child u1 {
    .clk,
    "data_(.*)" => "bus_\1"
};
\end{lstlisting}

This maps \code{data\_0} $\to$ \code{bus\_0}, \code{data\_1} $\to$ \code{bus\_1}, etc. Signal widths are inferred from the target module's port declarations.

Regex connections are processed in order and can be mixed with explicit connections. The priority order is:
\begin{enumerate}
  \item Explicit connections (highest priority)
  \item Same-name shorthand
  \item Regex patterns (first matching rule wins)
\end{enumerate}

\subsection{Regex Functions}

Use \code{\$func(...)} syntax to transform captured groups in replacements:

\begin{lstlisting}[language=Slip]
child u1 {
    .clk,
    "data_(.*)" => "bus_\$upper(\\1)",    // uppercase
    "ch_(.*)"   => "port_\$add(\\1, 1)"   // arithmetic
};
\end{lstlisting}

Built-in functions:
\begin{itemize}
  \item \textbf{String}: \code{\$upper}, \code{\$lower}, \code{\$reverse}, \code{\$substr}, \code{\$replace}, \code{\$concat}
  \item \textbf{Arithmetic}: \code{\$add}, \code{\$sub}, \code{\$mul}, \code{\$div}, \code{\$mod}
  \item \textbf{Bit}: \code{\$bit\_reverse}, \code{\$bit\_select}
\end{itemize}

Multiple functions and prefix/suffix text are supported:
\begin{lstlisting}[language=Slip]
"port_(.*)" => "pre_\$upper(\\1)_post"   // prefix + function + suffix
"a_(.*)_b_(.*)" => "\$upper(\\1)_\$lower(\\2)"  // multiple functions
\end{lstlisting}

\subsection{Custom Functions}

Register custom functions via the Python API:

\begin{lstlisting}[language=Python]
from slip.semantic.regex_funcs import register

def my_func(s: str) -> str:
    return f"prefix_{s}"

register("my_func", my_func)
\end{lstlisting}

Then use in Slip:

\begin{lstlisting}[language=Slip]
child u1 { .clk, "data_(.*)" => "\$my_func(\\1)" };
\end{lstlisting}

Custom functions can override built-in functions. The function signature must accept string arguments and return a string.

\subsection{Native Custom Functions (defun)}

Define custom functions directly in Slip using Python-style syntax:

\begin{lstlisting}[language=Slip]
defun prefix(s):
    return "bus_" + s

defun to_upper(s):
    return s.upper()

defun inc(n):
    return n + 1

module top (clk) {
    logic clk;
    inner u1 {
        .clk,
        "data_(.*)" => "\$prefix(\\1)",
        "name_(.*)" => "\$to_upper(\\1)",
        "ch_(.*)" => "\$inc(\\1)"
    };
}
\end{lstlisting}

The function body supports full Python expressions, including:
\begin{itemize}
  \item String concatenation: \code{a + b}
  \item String methods: \code{.upper()}, \code{.lower()}, \code{.reverse()}, \code{.replace()}, \code{.strip()}, \code{.split()}, \code{.startswith()}, etc.
  \item Arithmetic: \code{+}, \code{-}, \code{*}, \code{/}, \code{\%}
  \item Method chaining: \code{s.upper().reverse()}
  \item Built-in functions: \code{len()}, \code{int()}, \code{str()}, \code{range()}, etc.
  \item Indexing and slicing: \code{s[0]}, \code{s[1:3]}, \code{s[3:]}, \code{s[:3]}, \code{s[-1]}
  \item Ternary expressions: \code{x if cond else y}
\end{itemize}

\subsubsection{Examples}

Replace characters in port names:
\begin{lstlisting}[language=Slip]
defun fix_names(s):
    return s.replace("_", "-")

module inner (data_in, data_out) {
    logic data_in;
    logic data_out;
    assign data_out = data_in;
}

module top (clk) {
    logic clk;
    inner u1 { .clk, "data_(.*)" => "bus_\$fix_names(\\1)" };
    // data_in -> bus_data-in, data_out -> bus_data-out
}
\end{lstlisting}

Extract part of port name using slicing:
\begin{lstlisting}[language=Slip]
defun shorten(s):
    return s[:4]  // take first 4 characters

module inner (channel_a, channel_b) {
    logic channel_a;
    logic channel_b;
    assign channel_b = channel_a;
}

module top (clk) {
    logic clk;
    inner u1 { .clk, "channel_(.*)" => "\$shorten(\\1)" };
    // channel_a -> chan, channel_b -> chan
}
\end{lstlisting}

Combine methods with chaining:
\begin{lstlisting}[language=Slip]
defun clean_upper(s):
    return s.strip().upper().replace("_", "")

module inner (sig_in, sig_out) {
    logic sig_in;
    logic sig_out;
    assign sig_out = sig_in;
}

module top (clk) {
    logic clk;
    inner u1 { .clk, "sig_(.*)" => "port_\$clean_upper(\\1)" };
    // sig_in -> port_SIGIN, sig_out -> port_SIGOUT
}
\end{lstlisting}

Conditional logic with \code{startswith()}:
\begin{lstlisting}[language=Slip]
defun prefix_by_type(s):
    return "axi_" + s if s.startswith("ar") else "apb_" + s

module inner (araddr, awaddr, wdata) {
    logic araddr;
    logic awaddr;
    logic wdata;
    assign wdata = araddr;
}

module top (clk) {
    logic clk;
    inner u1 { .clk, "(.*)" => "\$prefix_by_type(\\1)" };
    // araddr -> axi_araddr, awaddr -> apb_awaddr, wdata -> apb_wdata
}
\end{lstlisting}

Functions can be defined anywhere at the top level (before or after modules). All \code{defun} functions are registered before instance resolution.

\section{Port Connection Constants}

Port connection expressions that are constants must follow these rules:

\begin{itemize}
  \item \textbf{Allowed}: width-prefixed literals (e.g., \code{1'b0}, \code{8'hFF}) and fill constants (\code{'0}, \code{'1}).
  \item \textbf{Forbidden}: bare integers (e.g., \code{0}, \code{42}).
\end{itemize}

\begin{lstlisting}[language=Slip]
// WRONG -- bare integer in port connection
inner u1 { .data_in(0) };

// OK -- width-prefixed
inner u1 { .data_in(1'b0) };

// OK -- fill constant (width determined by target port)
inner u1 { .data_in('0) };
\end{lstlisting}

The fill constants \code{'0} and \code{'1} are SystemVerilog native adaptive-width constants. Slip passes them through directly; the downstream tool determines their width from the target port.

% ──────────────────────────────────────────────────────────────────
\chapter{Metaprogramming}
% ──────────────────────────────────────────────────────────────────

Slip provides compile-time \slipinline{`for} loops and \slipinline{`if} conditionals. These are fully expanded before code generation.

\section{\texorpdfstring{\code{`for}}{`for} Loop}

Syntax uses backtick-prefixed keywords and loop variables:

\begin{lstlisting}[language=Slip]
`for (`i = 0; `i < 4; `i = `i + 1) {
    logic [7:0] data_`i;
    assign data_`i = `i;
}
\end{lstlisting}

After expansion, this becomes:

\begin{lstlisting}[language=Slip]
logic [7:0] data_0;
logic [7:0] data_1;
logic [7:0] data_2;
logic [7:0] data_3;
assign data_0 = 0;
assign data_1 = 1;
assign data_2 = 2;
assign data_3 = 3;
\end{lstlisting}

\subsection{Template Identifiers}

Loop variables can appear inside identifier names using backtick syntax. The lexer merges the surrounding identifier text with the backtick variable into a single token:

\begin{lstlisting}[language=Slip]
`for (`i = 0; `i < 3; `i = `i + 1) {
    assign out_`i = in_`i;
}
\end{lstlisting}

Generates:

\begin{lstlisting}[language=SystemVerilog]
    logic out_0;
    logic in_0;
    logic out_1;
    logic in_1;
    logic out_2;
    logic in_2;

    assign out_0 = in_0;
    assign out_1 = in_1;
    assign out_2 = in_2;
\end{lstlisting}

\subsection{Nested Loops}

\begin{lstlisting}[language=Slip]
`for (`i = 0; `i < 2; `i = `i + 1) {
    `for (`j = 0; `j < 2; `j = `j + 1) {
        assign result_`i_`j = `i + `j;
    }
}
\end{lstlisting}

Generates:

\begin{lstlisting}[language=SystemVerilog]
    logic result_0_0;
    logic result_0_1;
    logic result_1_0;
    logic result_1_1;

    assign result_0_0 = 0;
    assign result_0_1 = 1;
    assign result_1_0 = 1;
    assign result_1_1 = 2;
\end{lstlisting}

Template identifiers with multiple backtick variables (e.g., \code{a`i\_`j}) are fully supported. The lexer merges them into a single token in a single pass.

\subsection{Constant Folding}

\label{sec:constant-folding}

After substituting the loop variable with its current integer value, the expansion engine performs \textbf{constant folding} on the resulting expression tree:

\begin{table}[htbp]
\centering
\begin{tabular}{ll}
\toprule
\textbf{Pattern} & \textbf{Result} \\
\midrule
Both operands are integer literals & Evaluate to result literal \\
One operand is \code{0}, op is \code{+} or \code{-} & Identity: return other operand \\
One operand is \code{1}, op is \code{*} or \code{/} & Identity: return other operand \\
One operand is \code{0}, op is \code{*} & Zero: return \code{0} \\
Unary operand is integer literal & Evaluate to result literal \\
Ternary condition is integer literal & Select corresponding branch \\
Parenthesized expression is literal & Unwrap parentheses \\
\bottomrule
\end{tabular}
\end{table}

Subexpressions containing non-constant references (e.g., parameters) are preserved unchanged. Only the constant parts are simplified.

\textbf{Example} --- bit reversal with parameterized width:

\begin{lstlisting}[language=Slip]
module bit_reverse #(param W = 8) (din, dout) {
    logic [W-1:0] din;
    logic [W-1:0] dout;

    comb {
        `for (`i = 0; `i < W; `i = `i + 1) {
            dout[W - 1 - `i] = din[`i];
        }
    }
}
\end{lstlisting}

When \code{W=8} and \code{`i=0}, the index expression \code{W - 1 - 0} is folded:
\begin{enumerate}
  \item Substitute \code{`i} with \code{0}: \code{W - 1 - 0}
  \item Fold \code{1 - 0} $\to$ \code{1}: \code{W - 1}
  \item Result: \code{dout[W - 1] = din[0]}
\end{enumerate}

The generated SystemVerilog contains no \code{- 0} artifacts.

\subsection{Compile-Time Expression Evaluation}

\slipinline{`for} init, bound, and step expressions, as well as \slipinline{`if} conditions, must be evaluable to integer constants at compile time. The evaluation engine supports:

\begin{itemize}
  \item Integer literals and \code{param}/\code{localparam} references
  \item Arithmetic: \code{+}, \code{-}, \code{*}, \code{/}, \code{\%}, \code{**}
  \item Relational: \code{<}, \code{>}, \code{<=}, \code{>=}, \code{==}, \code{!=}, \code{===}, \code{!==}
  \item Logical: \code{\&\&}, \code{||}
  \item Bitwise: \code{\&}, \code{|}, \code{\^}, \code{\~}, \code{<<}, \code{>>}, \code{<<<}, \code{>>>}
  \item Unary prefix: \code{-}, \code{!}, \code{\~}
  \item Parenthesized sub-expressions
\end{itemize}

\section{\texorpdfstring{\code{`if}}{`if} Conditional}

\begin{lstlisting}[language=Slip]
`if (ENABLE_DEBUG) {
    logic [7:0] debug_out;
    assign debug_out = internal_signal;
}
`else {
    assign debug_out = 0;
}
\end{lstlisting}

When \code{ENABLE\_DEBUG} evaluates to non-zero, generates:

\begin{lstlisting}[language=SystemVerilog]
    logic [7:0] debug_out;
    assign debug_out = internal_signal;
\end{lstlisting}

When \code{ENABLE\_DEBUG} evaluates to zero, generates:

\begin{lstlisting}[language=SystemVerilog]
assign debug_out = 0;
\end{lstlisting}

The condition must be a compile-time constant expression (parameter, localparam, or integer literal expression).

\section{Restrictions}

\begin{itemize}
  \item Loop variables (\slipinline{`i}) cannot be assigned within the loop body --- they are controlled by the expansion engine.
  \item \slipinline{`for} / \slipinline{`if} conditions must be evaluable at compile time.
  \item The step variable must match the loop variable (e.g., \slipinline{`i = `i + 1}).
  \item Maximum 1024 iterations per loop (prevents infinite loops).
\end{itemize}

% ──────────────────────────────────────────────────────────────────
\chapter{Expressions and Operators}
% ──────────────────────────────────────────────────────────────────

\section{Operator Precedence}

From highest to lowest (following IEEE 1800-2017):

\begin{table}[htbp]
\centering
\begin{longtable}{cll}
\toprule
\textbf{Priority} & \textbf{Operators} & \textbf{Associativity} \\
\midrule
\endfirsthead
\toprule
\textbf{Priority} & \textbf{Operators} & \textbf{Associativity} \\
\midrule
\endhead
1 (highest) & \code{[]}, \code{[ : ]}, function call, \code{N'(expr)} & Left \\
2 & \code{+}, \code{-}, \code{!}, \code{\~} (unary); \code{signed'}, \code{unsigned'} & Right \\
3 & \code{**} & Right \\
4 & \code{*}, \code{/}, \code{\%} & Left \\
5 & \code{+}, \code{-} (binary) & Left \\
6 & \code{<<}, \code{>>}, \code{<<<}, \code{>>>} & Left \\
7 & \code{<}, \code{>}, \code{<=}, \code{>=} & Left \\
8 & \code{==}, \code{!=}, \code{===}, \code{!==} & Left \\
9 & \code{\&} & Left \\
10 & \code{\^} & Left \\
11 & \code{|} & Left \\
12 & \code{\&\&} & Left \\
13 & \code{||} & Left \\
14 & \code{?:} (ternary) & Right \\
15 (lowest) & \code{inside} (set membership) & Left \\
\bottomrule
\end{longtable}
\end{table}

\section{Bit Selection and Part Select}

\begin{lstlisting}[language=Slip]
data[3]          // bit select
data[7:0]        // part select
data[i]          // variable index
\end{lstlisting}

\section{Concatenation and Replication}

\begin{lstlisting}[language=Slip]
{a, b, c}        // concatenation
{4{1'b1}}        // replication: 4 copies of 1'b1
\end{lstlisting}

Nested replications are supported: \code{\{4\{\{2\{1'b0\}\}\}\}} produces 8 bits of zero.

\section{Type Casting}

\begin{lstlisting}[language=Slip]
signed'(expr)    // signed cast
unsigned'(expr)  // unsigned cast
8'(expr)         // width cast (truncate or zero-extend to 8 bits)
\end{lstlisting}

Both \code{signed} and \code{unsigned} without a following \code{'} raise a syntax error.

\section{Ternary Operator}

\begin{lstlisting}[language=Slip]
sel ? a : b
\end{lstlisting}

Right-associative: \code{a ? b : c ? d : e} is parsed as \code{a ? b : (c ? d : e)}.

% ──────────────────────────────────────────────────────────────────
\chapter{Special Constants}
% ──────────────────────────────────────────────────────────────────

\section{Fill Constants}

\begin{lstlisting}[language=Slip]
'0    // all bits zero (width determined by context)
'1    // all bits one  (width determined by context)
\end{lstlisting}

These are SystemVerilog native fill constants. Their width is automatically determined by the connected port or assignment target. Slip passes them through without width inference.

\section{Width-Prefixed Literals}

\begin{lstlisting}[language=Slip]
1'b0          // 1-bit binary zero
1'b1          // 1-bit binary one
8'hFF         // 8-bit hexadecimal
4'b1010       // 4-bit binary
16'd255       // 16-bit decimal
8'o377        // 8-bit octal
\end{lstlisting}

Digits are validated against the specified radix. For example, \code{4'bABCD} and \code{4'o888} are rejected. The four-state values \code{x} and \code{z} are accepted in any radix.

\section{Bare Integers}

\begin{lstlisting}[language=Slip]
0             // bare integer (OK in assignments and expressions)
42            // bare integer
\end{lstlisting}

Bare integers are \textbf{not allowed} in port connections --- use width-prefixed literals or fill constants instead.

% ════════════════════════════════════════════════════════════════════
% Part III: Advanced Topics
% ════════════════════════════════════════════════════════════════════

\chapter{IP Integration}

Slip can instantiate external SystemVerilog modules by providing IP directories or VCS-format filelists:

\begin{lstlisting}[language=bash]
slip build -ip ./ip_cores -ip ./vendor_lib design.slip
slip build -f ip.f design.slip
slip build -f ip.f -ip ./ip_cores design.slip
\end{lstlisting}

\section{How It Works}

\begin{enumerate}
  \item Slip collects IP sources from filelists (\code{-f}) and IP directories (\code{-ip}).
  \item IP directories are scanned recursively for \code{.sv} and \code{.v} files.
  \item All IP files are parsed together with a shared preprocessor context, so \code{`define} macros propagate across files.
  \item When an instance references a module not defined in the \code{.slip} source, Slip uses pyslang to reflect the external module's ports and parameters.
  \item Port widths are evaluated using the instance's parameter overrides.
  \item Signals are created with correct widths for regex-mapped connections.
\end{enumerate}

\section{Filelist Support}

Slip supports VCS-format filelist (\code{.f}) files for explicit control over IP source ordering, include paths, and macro definitions:

\begin{lstlisting}[language=Slip]
// ip.f
+incdir+./include
+define+DATA_WIDTH=32
+define+ADDR_WIDTH=8
./rtl/fifo.sv
./rtl/arbiter.sv
-f vendor_lib.f
\end{lstlisting}

\begin{itemize}
  \item \code{+incdir+path} --- include search directory (multiple paths separated by \code{+})
  \item \code{+define+MACRO=value} --- macro definition
  \item \code{-f sub.f} --- nested filelist reference (recursive, with circular reference detection)
  \item File paths are resolved relative to the filelist location
  \item Duplicate files are automatically deduplicated
  \item Comments (\code{//}) and empty lines are ignored
\end{itemize}

\section{Two-Pass Scanning}

IP scanning uses a two-pass approach to resolve \code{`define} ordering:

\textbf{Pass 1 --- Collection}: Parse all filelists, collect source files, include paths, and macro definitions. Recursively scan IP directories and merge with filelist entries (deduplication).

\textbf{Pass 2 --- Reflection}: Parse all IP files together using \code{SyntaxTree.fromFiles()} with a shared preprocessor context. This ensures that \code{`define} macros in one file are available to subsequent files, solving the define ordering problem that would occur with per-file parsing.

\section{Example}

Given an external IP file \code{ip\_cores/fifo.sv}:

\begin{lstlisting}[language=SystemVerilog]
module fifo #(
    parameter DEPTH = 16,
    parameter WIDTH = 8
)(
    input  logic             clk,
    input  logic             rst_n,
    input  logic [WIDTH-1:0] din,
    output logic [WIDTH-1:0] dout,
    output logic             full
);
\end{lstlisting}

In your Slip source:

\begin{lstlisting}[language=Slip]
module top (clk, rst_n, data_in, data_out) {
    logic clk;
    logic rst_n;
    logic [7:0] data_in;
    logic [7:0] data_out;

    fifo #(.DEPTH(32)) u_fifo {
        .clk,
        .rst_n,
        .din(data_in),
        .dout(data_out),
        .full(_)
    };
}
\end{lstlisting}

The \code{.full(\_)} connection leaves the \code{full} output dangling.

% ──────────────────────────────────────────────────────────────────
\chapter{Semantic Rules}
% ──────────────────────────────────────────────────────────────────

\section{Multi-Driver Detection}

A signal cannot be driven by multiple procedural blocks. This is caught at compile time:

\begin{lstlisting}[language=Slip]
// ERROR: signal 'q' driven by multiple blocks
seq (clk) { q <= a; }
seq (clk) { q <= b; }
\end{lstlisting}

The analyzer tracks assignments through \code{if} and \code{for} nesting within each block. Two separate \code{seq} or \code{comb} blocks driving the same signal produce an error.

\section{Combinational Loop Detection}

Feedback loops inside \code{comb} blocks are detected using DFS-based cycle detection on the signal dependency graph:

\begin{lstlisting}[language=Slip]
// ERROR: combinational loop detected
comb {
    a = b + 1;
    b = a + 1;  // 'a' depends on 'b', 'b' depends on 'a'
}
\end{lstlisting}

The detector builds a dependency graph per \code{comb} block by tracing which signals are read and which are driven. If a cycle is found, a \code{SlipSemanticError} is raised.

\section{Width Mismatch Warnings}

The compiler checks for width mismatches in assignments and instance port connections. When both sides have concrete integer widths and they differ, a warning is emitted:

\begin{lstlisting}[language=Slip]
module m (out[7:0]) {
    assign out = 42;  // WARNING: 'out' is 8 bits, but value is 32 bits
}
\end{lstlisting}

Width-prefixed literals use their declared width:

\begin{lstlisting}[language=Slip]
module m (out[7:0]) {
    assign out = 8'hFF;  // OK: both 8 bits
}
\end{lstlisting}

Instance port connections are also checked:

\begin{lstlisting}[language=Slip]
module sub (data[7:0]) { ... }
module top (out[3:0]) {
    sub u { .data(out) };  // WARNING: port expects 8 bits, signal is 4 bits
}
\end{lstlisting}

Parameterized widths with default values are evaluated using those defaults. Widths that cannot be evaluated are silently skipped.

\section{Localparam Override Protection}

Parameters can be overridden at instantiation; localparams cannot:

\begin{lstlisting}[language=Slip]
module example #(param W = 8) () {
    localparam MASK = 255;
}

// OK
example #(.W(16)) u1 {};

// ERROR: cannot override localparam 'MASK'
example #(.MASK(511)) u2 {};
\end{lstlisting}

\section{Implicit Signal Width Inference}

Implicit signals connected to instance ports automatically inherit the target port's width:

\begin{lstlisting}[language=Slip]
module top (a, b) {
    logic [7:0] a;
    logic [7:0] b;
    // 'sum' is implicit --- inherits [7:0] from adder.sum port
    adder u1 { .a, .b, .sum };
}
\end{lstlisting}

Generates:

\begin{lstlisting}[language=SystemVerilog]
    // width from adder.sum
    logic [7:0] sum;
\end{lstlisting}

Other implicit signals (not connected to ports) default to 1-bit and are annotated with \code{// implicit, no width} in the generated output.

\section{Generated Code Annotations}

Slip adds review annotations to generated SystemVerilog signals and ports:

\begin{table}[htbp]
\centering
\begin{tabular}{ll}
\toprule
\textbf{Annotation} & \textbf{Meaning} \\
\midrule
\code{// implicit} & Signal or port was auto-generated (not explicitly declared) \\
\code{// no width} & Signal or port has no declared width (defaults to 1-bit) \\
\code{// width from X.Y} & Width was automatically inferred from instance port \\
\bottomrule
\end{tabular}
\end{table}

These annotations help identify signals that may need explicit width declarations during code review. Use \code{grep "no width" *.sv} to quickly find all un-sized signals.

\section{Integer Literal Validation}

Verilog-style sized literals are validated at lex time:
\begin{itemize}
  \item Digits must be valid for the specified radix (e.g., \code{4'bABCD} is rejected).
  \item The digit portion must not be empty or underscore-only (e.g., \code{4'b\_} is rejected).
  \item Four-state values \code{x}/\code{z} are accepted in any radix.
\end{itemize}

% ──────────────────────────────────────────────────────────────────
\chapter{Error Reference}
% ──────────────────────────────────────────────────────────────────

\section{Compilation Errors}

\begin{description}[style=nextline,leftmargin=2em]

\item[\texttt{cannot override localparam 'name' of module 'M' (instance 'I')}]
Attempting to override a localparam in instance \code{\#(...)}.

\item[\texttt{unconnected input port 'name' in instance 'I' of module 'M'}]
Input port not connected and not marked with \code{\_}. The error includes active regex rules when applicable.

\item[\texttt{unable to infer width for implicit signal 'name'}]
Implicit signal used in a context where width cannot be determined (legacy; now defaults to 1-bit with \code{// implicit, no width} annotation).

\item[\texttt{cyclic module dependency detected among: A, B}]
Modules have circular instantiation dependency (e.g., module A instantiates B, and B instantiates A). This is illegal in synthesizable hardware.

\item[\texttt{bare integer '...' in port connection must have explicit width}]
Bare integer used in port connection. Use a width-prefixed literal (\texttt{1'b0}) or a fill constant (\texttt{'0}).

\item[\texttt{signal 'name' driven by multiple blocks}]
Multi-driver conflict --- the signal is driven by more than one \code{seq} or \code{comb} block.

\item[\texttt{invalid digit 'X' for radix 'r' in literal}]
Integer literal contains a digit invalid for the specified radix (e.g.\ \texttt{4'bABCD}).

\item[\texttt{integer literal has no actual digit value}]
Literal has only underscores after the radix prefix (e.g.\ \texttt{4'b\_}).

\item[\texttt{`for loop exceeded 1024 iterations}]
Metaprogramming loop body was expanded more than 1024 times.

\item[\texttt{cannot evaluate 'name' at compile time}]
Non-constant expression in \slipinline{`for}/\slipinline{`if} condition, init, or step.

\item[\texttt{`for step variable '...' must match loop variable '...'}]
Step variable differs from the declared loop variable. The step expression determines the next value of the loop variable, so the step variable must match the loop variable.

\item[\texttt{unsupported operator '...' in compile-time expression}]
Operator not supported by the constant expression evaluator.

\item[\texttt{combinational loop detected: ...}]
A feedback loop was found inside a \code{comb} block. The error message lists the signals involved in the cycle.

\item[\texttt{duplicate declaration of 'name'}]
A signal, instance, or localparam with the same name is declared more than once in the same module.

\item[\texttt{reset polarity inconsistency: ...}]
A reset signal is used with conflicting polarity (both \code{pos:} and \code{neg:}) across different \code{seq} blocks.

\end{description}

\section{Warnings}

\begin{description}[style=nextline,leftmargin=2em]

\item[\texttt{blocking '=' converted to '<=' in seq block}]
Automatic assignment correction in sequential logic.

\item[\texttt{pyslang not installed --- skipping validation}]
Install \code{pyslang} to enable IEEE 1800 validation of generated output.

\item[\texttt{width mismatch: 'name' is N bits, but value is M bits}]
An assignment has mismatched widths between the target signal and the value expression.

\item[\texttt{port width mismatch: 'port' expects N bits, but 'signal' is M bits}]
An instance port connection has mismatched widths between the port declaration and the connected signal.

\end{description}

% ──────────────────────────────────────────────────────────────────
\chapter{Complete Examples}
% ──────────────────────────────────────────────────────────────────

\section{Parameterized Pipeline Register}

\begin{lstlisting}[language=Slip]
module pipeline_reg #(param WIDTH = 8) (clk, rst_n, d, q) {
    logic clk;
    logic rst_n;
    logic [WIDTH-1:0] d;
    logic [WIDTH-1:0] q;

    seq (clk, neg: rst_n) {
        if (!rst_n) {
            q <= 0;
        } else {
            q <= d;
        }
    }
}
\end{lstlisting}

Generates:

\begin{lstlisting}[language=SystemVerilog]
module pipeline_reg #(
    parameter WIDTH = 8
) (
    input logic clk,
    input logic rst_n,
    input logic [WIDTH - 1:0] d,
    output logic [WIDTH - 1:0] q
);

    always_ff @(posedge clk or negedge rst_n) begin
        if (!rst_n) begin
            q <= 0;
        end else begin
            q <= d;
        end
    end

endmodule
\end{lstlisting}

\section{Bit Reversal with Metaprogramming}

\begin{lstlisting}[language=Slip]
module bit_reverse #(param W = 8) (din, dout) {
    logic [W-1:0] din;
    logic [W-1:0] dout;

    comb {
        `for (`i = 0; `i < W; `i = `i + 1) {
            dout[W - 1 - `i] = din[`i];
        }
    }
}
\end{lstlisting}

Generates:

\begin{lstlisting}[language=SystemVerilog]
module bit_reverse #(
    parameter W = 8
) (
    input logic [W - 1:0] din,
    output logic [W - 1:0] dout
);

    always_comb begin
        dout[W - 1] = din[0];
        dout[W - 1 - 1] = din[1];
        dout[W - 1 - 2] = din[2];
        dout[W - 1 - 3] = din[3];
        dout[W - 1 - 4] = din[4];
        dout[W - 1 - 5] = din[5];
        dout[W - 1 - 6] = din[6];
        dout[W - 1 - 7] = din[7];
    end

endmodule
\end{lstlisting}

Note: constant folding simplifies \code{W - 1 - 0} to \code{W - 1}, but leaves expressions like \code{W - 1 - 1} unfolded because they contain non-constant subexpressions. The downstream synthesis tool evaluates these at elaboration time.

\section{State Machine with Localparams}

\begin{lstlisting}[language=Slip]
module arbiter (clk, rst_n, req, grant) {
    logic clk;
    logic rst_n;
    logic [3:0] req;
    logic [3:0] grant;
    logic [1:0] cur;

    localparam S_IDLE = 0;
    localparam S_GNT0 = 1;
    localparam S_GNT1 = 2;

    seq (clk, neg: rst_n) {
        if (!rst_n) {
            cur <= S_IDLE;
        } else {
            case (cur) {
                S_IDLE: {
                    if (req[0]) cur <= S_GNT0;
                    if (req[1]) cur <= S_GNT1;
                }
                S_GNT0: {
                    if (!req[0]) cur <= S_IDLE;
                }
                S_GNT1: {
                    if (!req[1]) cur <= S_IDLE;
                }
                default: cur <= S_IDLE;
            }
        }
    }

    assign grant[0] = (cur == S_GNT0);
    assign grant[1] = (cur == S_GNT1);
}
\end{lstlisting}

Generates:

\begin{lstlisting}[language=SystemVerilog]
module arbiter (
    input logic clk,
    input logic rst_n,
    input logic [3:0] req,
    output logic [3:0] grant
);

    localparam S_IDLE = 0;
    localparam S_GNT0 = 1;
    localparam S_GNT1 = 2;

    logic [1:0] cur;

    assign grant[0] = (cur == S_GNT0);
    assign grant[1] = (cur == S_GNT1);

    always_ff @(posedge clk or negedge rst_n) begin
        if (!rst_n) begin
            cur <= S_IDLE;
        end else begin
            case (cur)
                S_IDLE: begin
                    if (req[0]) begin
                        cur <= S_GNT0;
                    end
                    if (req[1]) begin
                        cur <= S_GNT1;
                    end
                end
                S_GNT0: begin
                    if (!req[0]) begin
                        cur <= S_IDLE;
                    end
                end
                S_GNT1: begin
                    if (!req[1]) begin
                        cur <= S_IDLE;
                    end
                end
                default: begin
                    cur <= S_IDLE;
                end
            endcase
        end
    end

endmodule
\end{lstlisting}

\section{Arbiter with Case Statement}

\begin{lstlisting}[language=Slip]
module arbiter_case (clk, rst_n, req, grant) {
    logic clk;
    logic rst_n;
    logic [3:0] req;
    logic [3:0] grant;
    logic [2:0] cur;

    localparam S_IDLE = 0;
    localparam S_GNT0 = 1;
    localparam S_GNT1 = 2;
    localparam S_GNT2 = 3;
    localparam S_GNT3 = 4;

    seq (clk, neg: rst_n) {
        if (!rst_n) {
            cur <= S_IDLE;
        } else {
            case (cur) {
                S_IDLE: {
                    if (req[0]) cur <= S_GNT0;
                    if (req[1]) cur <= S_GNT1;
                    if (req[2]) cur <= S_GNT2;
                    if (req[3]) cur <= S_GNT3;
                }
                S_GNT0: {
                    if (!req[0]) cur <= S_IDLE;
                }
                S_GNT1: {
                    if (!req[1]) cur <= S_IDLE;
                }
                default: cur <= S_IDLE;
            }
        }
    }

    comb {
        grant = 0;
        case (cur) {
            S_GNT0: grant[0] = 1;
            S_GNT1: grant[1] = 1;
            S_GNT2: grant[2] = 1;
            S_GNT3: grant[3] = 1;
            default: grant = 0;
        }
    }
}
\end{lstlisting}

Generates:

\begin{lstlisting}[language=SystemVerilog]
module arbiter_case (
    input logic clk,
    input logic rst_n,
    input logic [3:0] req,
    output logic [3:0] grant
);

    localparam S_IDLE = 0;
    localparam S_GNT0 = 1;
    localparam S_GNT1 = 2;
    localparam S_GNT2 = 3;
    localparam S_GNT3 = 4;

    logic [2:0] cur;

    always_ff @(posedge clk or negedge rst_n) begin
        if (!rst_n) begin
            cur <= S_IDLE;
        end else begin
            case (cur)
                S_IDLE: begin
                    if (req[0]) begin
                        cur <= S_GNT0;
                    end
                    if (req[1]) begin
                        cur <= S_GNT1;
                    end
                    if (req[2]) begin
                        cur <= S_GNT2;
                    end
                    if (req[3]) begin
                        cur <= S_GNT3;
                    end
                end
                S_GNT0: begin
                    if (!req[0]) begin
                        cur <= S_IDLE;
                    end
                end
                S_GNT1: begin
                    if (!req[1]) begin
                        cur <= S_IDLE;
                    end
                end
                default: begin
                    cur <= S_IDLE;
                end
            endcase
        end
    end

    always_comb begin
        grant = 0;
        case (cur)
            S_GNT0: begin
                grant[0] = 1;
            end
            S_GNT1: begin
                grant[1] = 1;
            end
            S_GNT2: begin
                grant[2] = 1;
            end
            S_GNT3: begin
                grant[3] = 1;
            end
            default: begin
                grant = 0;
            end
        endcase
    end

endmodule
\end{lstlisting}

\section{Hierarchical Design}

\begin{lstlisting}[language=Slip]
module stage #(param W = 8) (clk, rst_n, data_in, data_out) {
    logic clk;
    logic rst_n;
    logic [W-1:0] data_in;
    logic [W-1:0] data_out;
    logic [W-1:0] reg_data;

    seq (clk, neg: rst_n) {
        if (!rst_n) {
            reg_data <= 0;
        } else {
            reg_data <= data_in;
        }
    }

    assign data_out = reg_data;
}

module pipeline (clk, rst_n, din, dout) {
    logic clk;
    logic rst_n;
    logic [7:0] din;
    logic [7:0] dout;
    logic [7:0] stage_data [0:3];

    assign stage_data[0] = din;

    `for (`i = 1; `i < 3; `i = `i + 1) {
        stage #(.W(8)) u_stage_`i {
            .clk,
            .rst_n,
            .data_in(stage_data[`i - 1]),
            .data_out(stage_data[`i])
        };
    }

    assign dout = stage_data[2];
}
\end{lstlisting}

Generates:

\begin{lstlisting}[language=SystemVerilog]
module stage #(
    parameter W = 8
) (
    input logic clk,
    input logic rst_n,
    input logic [W - 1:0] data_in,
    output logic [W - 1:0] data_out
);

    logic [W - 1:0] reg_data;

    assign data_out = reg_data;

    always_ff @(posedge clk or negedge rst_n) begin
        if (!rst_n) begin
            reg_data <= 0;
        end else begin
            reg_data <= data_in;
        end
    end

endmodule

module pipeline (
    input logic clk,
    input logic rst_n,
    input logic [7:0] din,
    output logic [7:0] dout
);

    logic [7:0] stage_data [0:3];

    assign stage_data[0] = din;
    assign dout = stage_data[2];

    stage #(
        .W(8)
    ) u_stage_1 (
        .clk(clk),
        .rst_n(rst_n),
        .data_in(stage_data[0]),
        .data_out(stage_data[1])
    );

    stage #(
        .W(8)
    ) u_stage_2 (
        .clk(clk),
        .rst_n(rst_n),
        .data_in(stage_data[1]),
        .data_out(stage_data[2])
    );

endmodule
\end{lstlisting}

% ════════════════════════════════════════════════════════════════════
\appendix
% ════════════════════════════════════════════════════════════════════

\chapter{Formal Grammar (EBNF)}

\begin{lstlisting}[basicstyle=\ttfamily\footnotesize,
  keywordstyle={},
  frame=single,
  backgroundcolor=\color{codebg},
  escapechar=@]
CompilationUnit ::= { TopLevelItem }

TopLevelItem ::= ModuleDecl | FuncDef | IncludeDirective

IncludeDirective ::= "include" STRING ";"

ModuleDecl ::= "module" IDENT
               [ "#(" ParameterList ")" ]
               [ "(" PortList ")" ]
               "{" { Statement } "}"

ParameterList ::= ParameterDecl { "," ParameterDecl }
ParameterDecl ::= "param" IDENT "=" Expr

PortList      ::= PortItem { "," PortItem }
PortItem      ::= IDENT [ "[" Expr [ ":" Expr ] "]" ]

Statement ::= SignalDecl
            | LocalparamDecl
            | AssignStmt
            | SeqBlock
            | CombBlock
            | InitialBlock
            | IfStmt
            | ForStmt
            | CaseStmt
            | GenForStmt
            | GenIfStmt
            | InstanceStmt
            | BlockStmt

SignalDecl ::= ["logic"] ["signed"] [ "[" Expr "]" ]
               IDENT [ "[" Expr ":" Expr "]" ] ";"

LocalparamDecl ::= "localparam" IDENT "=" Expr ";"

AssignStmt ::= ["assign"] Lvalue "=" Expr ";"
             | ["assign"] Lvalue CompoundOp Expr ";"

CompoundOp ::= "+=" | "-=" | "*=" | "/=" | "%="
             | "&=" | "|=" | "^=" | "<<=" | ">>="
             | "<<<=" | ">>>="

SeqBlock     ::= "seq" "(" IDENT
                 [ "," ("pos" | "neg") ":" IDENT ] ")" BlockStmt
CombBlock    ::= "comb" BlockStmt
InitialBlock ::= "initial" BlockStmt

IfStmt  ::= "if" "(" Expr ")" BlockStmt
            [ "else" BlockStmt ]

ForStmt ::= "for" "(" IDENT "=" Expr ";" Expr ";"
                   ForStep ")" BlockStmt
ForStep ::= IDENT "=" Expr
          | IDENT CompoundOp Expr
          | IDENT "++"

CaseStmt ::= ("case" | "casez" | "casex") "(" Expr ")"
             "{" { CaseItem } "}"
CaseItem ::= PatternList ":" Statement
           | "default" ":" Statement
PatternList ::= Expr { "," Expr }

GenForStmt ::= "`for" "(" TICK_IDENT "=" Expr ";" Expr ";"
                          TICK_IDENT "=" Expr ")" BlockStmt
GenIfStmt  ::= "`if" "(" Expr ")" BlockStmt
                         [ "`else" BlockStmt ]

InstanceStmt ::= IDENT
                 [ "#(" NamedParam { "," NamedParam } ")" ]
                 IDENT "{" { Connection } "}" [";"]

NamedParam ::= "." IDENT "(" Expr ")"
Connection ::= "." IDENT [ "(" Expr ")" ]
             | STRING "=>" STRING

BlockStmt ::= "{" { Statement } "}"
Lvalue    ::= IDENT { "[" Expr "]" | "[" Expr ":" Expr "]" }
\end{lstlisting}

\textbf{Notes on the grammar:}
\begin{itemize}
  \item \code{TICK\_IDENT} denotes a backtick-prefixed identifier (e.g., \code{`i}).
  \item Template identifiers (e.g., \code{data\_`i}) are formed at the lexical level by merging adjacent \code{IDENT}, \code{TICK}, and \code{TICK\_IDENT} tokens.
  \item Expressions are parsed by a Pratt parser and support all operators listed in the operator precedence table.
  \item Special expression forms include concatenation \code{\{a, b\}}, replication \code{\{n\{expr\}\}}, width cast \code{N'(expr)}, and set membership \code{x inside \{v1, v2, ...\}}.
  \item Compound assignment operators are desugared at parse time into regular assignments.
\end{itemize}

\end{document}
